\section{Síntesis y ampliación}

\subsection{Tabla de síntesis de toda la clase}
\begin{table}[H]
\centering
\begin{tabular}{p{2.8cm}p{4.1cm}p{4.6cm}}
\toprule
\textbf{Tema} & \textbf{Idea central del conferencista} & \textbf{Implicación para la investigación en Agentes} \\
\midrule
Tendencia de apertura & El aumento de capacidad y la disminución del acceso ocurren simultáneamente & Los límites de la investigación quedan moldeados directamente por los permisos de acceso \\
Marco de tres niveles de acceso & La estratificación en API / Open Weights / Open Source es clara & Cada nivel corresponde a una capacidad distinta de experimentación verificable \\
Ingeniería de Agentes basados en API & El LLM funge como controlador de ejecución & El énfasis se desplaza hacia la memoria, la reflexión y la gestión de trayectorias \\
Benchmark y evidencia empírica & Contar con un sistema evaluable es condición previa para el progreso & Se pasa de una narrativa basada en demos a comparaciones reproducibles \\
Definición de apertura y gobernanza & Los límites terminológicos influyen en las políticas y la colaboración & «Descargable» no equivale a «reproducible» \\
Ruta a largo plazo & La capacidad técnica y la infraestructura coevolucionan & La ciencia abierta requiere la coordinación de estándares, cadenas de herramientas y cómputo \\
\bottomrule
\end{tabular}
\caption{Síntesis de los argumentos centrales de Lecture 03}
\end{table}

\subsection{Lecturas adicionales}
\begin{itemize}
    \item Bommasani et al., \textit{On the Opportunities and Risks of Foundation Models}, 2021.
    \item Liang et al., \textit{Holistic Evaluation of Language Models (HELM)}, 2022.
    \item Artículos de referencia relacionados con Agent evaluation and reproducibility (incluida la serie de trabajos de MLAgentBench).
    \item Discusiones en curso y actualizaciones de borradores de la OSI sobre la definición de apertura de los modelos de IA.
    \item Prácticas de ingeniería en marcos de trabajo con múltiples Agentes: comparación de diseño entre AutoGen, LangGraph y CrewAI.
\end{itemize}

\subsection{Idea de cierre}
\textit{``Access shapes research.''} Esta frase, aparentemente simple, en realidad define los límites reales de los investigadores en la era de los Agentes.  
Si se desea lograr un progreso científico sostenible en las direcciones de Foundation Model y Agent, la comunidad debe avanzar simultáneamente en la capacidad de los modelos, el acceso abierto y la infraestructura reproducible, en lugar de perseguir únicamente un SOTA puntual.

%% --- Fin del contenido --- %%

\end{document}
